% V1.2 candidate method section. This file is intentionally not included in
% the frozen headline manuscript until the independent Monte Carlo gate passes.
\section{Joint Fusion--Bundle Optimization}\label{sec:v12-method}

\subsection{Target--Fusion--Observation Bundles}
Let $\mathcal M$ and $\mathcal Q$ denote the UAV and target sets.  A bistatic
observation $e=(i,j)\in\mathcal E_q$ is produced by transmitter $i$ and
receiver $j$ for target $q$.  Target $q$ is assigned to one fusion UAV
$f\in\mathcal M$.  The observation is local if $j=f$; otherwise, it consumes
one inter-UAV report from $j$ to $f$.  For each target--fusion pair $(q,f)$,
we define a bundle $b\in\mathcal B_{qf}$ as a feasible subset of
$\mathcal E_q$.  Its scheduler-side predicted detection probability, remote
report load, receiver-$j$ load, and fusion-processing load are denoted by
$p_{qfb}$, $r_{qfb}$, $a_{qfbj}$, and $c_{qfb}=|b|$, respectively.

The physical and detector models are unchanged from the true-erasure V1.1
system, but we make the RCS dependence explicit here.  For scheduler-side
target RCS $\sigma_q$, wavelength $\lambda$, and bistatic ranges $d_{iq}$ and
$d_{jq}$, the propagation gain and sensing SINR are
\begin{align}
g^{\rm s}_{ijq}(\sigma_q)
  &=\frac{\lambda^2\sigma_q}
  {(4\pi)^3d_{iq}^2d_{jq}^2},\\
\gamma^{\rm s}_{ijq}(\sigma_q)
  &=\frac{\widetilde P_i^{\rm s}G^{\rm hw}_{ij}G_p
  \eta^{\rm DD}_{ijq}\eta^{\rm col}_{ijq}\eta^{\rm wf}_{ijq}
  g^{\rm s}_{ijq}(\sigma_q)}
  {(N_0+I_j^{\rm s})(1+{\rm INR}^{\rm wf}_{ijq})}.
\label{eq:v12-rcs-sinr}
\end{align}
Here $G^{\rm hw}_{ij}=G_i^{\rm tx}G_j^{\rm rx}/L_{ij}^{\rm sys}$ is kept
separate from RCS.  The scheduler uses the declared mean RCS
$\sigma_q=\bar\sigma_q$ rather than a current-CPI realization.  Hence RCS is
an exogenous physical parameter that changes the precomputed utility
$p_{qfb}(\sigma_q)$; it is not a decision variable of the bundle master.

In particular, a failed remote packet contributes an exact zero soft
statistic.  All compared methods use the same finite-look mixture-calibrated
threshold in final Monte Carlo evaluation.  The bundle utility $p_{qfb}$ is a
moment-matched scheduling surrogate; it is not identified with empirical
$P_D$.

\subsection{Lexicographic Master Problem}
Let $u_{qfb}\in\{0,1\}$ indicate the selected bundle and let $d_q\geq0$ be
the deficit from the required detection probability $P_D^{\rm req}$.  We solve
\begin{equation}
\operatorname{lexmin}\left(
d_{\max},\ \sum_q d_q,\ \sum_{qfb}r_{qfb}u_{qfb},\
\sum_{qfb}c_{qfb}u_{qfb}\right),
\label{eq:v12-lex-objective}
\end{equation}
subject to
\begin{align}
\sum_{f,b}u_{qfb}&=1, &&\forall q,\label{eq:v12-one-bundle}\\
d_q&\geq P_D^{\rm req}-\sum_{f,b}p_{qfb}u_{qfb}, &&\forall q,\\
d_{\max}&\geq d_q, &&\forall q,\\
\sum_{q,b}u_{qfb}&\leq C_f^{\rm target}, &&\forall f,\\
\sum_{q,f,b}a_{qfbj}u_{qfb}&\leq C_j^{\rm rx}, &&\forall j,\\
\sum_{q,b}c_{qfb}u_{qfb}&\leq C_f^{\rm fusion}, &&\forall f,\\
\sum_{q,f,b}r_{qfb}u_{qfb}&\leq K_{\rm report},\\
\sum_{q,f,b}c_{qfb}u_{qfb}&\leq K_{\rm total}.
\end{align}
Per-target total- and local-observation limits are enforced during bundle
generation.  The empty bundle is retained because target assignment and
sensing-resource expenditure are distinct decisions.

When $\mathcal B_{qf}$ contains every feasible observation subset, this
formulation is equivalent to the original joint assignment--selection problem:
each original solution maps to one bundle per target, and each selected bundle
uniquely recovers the fusion destination and observation variables while
preserving every resource load.  We solve the four objectives in
\eqref{eq:v12-lex-objective} sequentially and lock each attained optimum before
optimizing the next tier.  The resulting integer solution is therefore
lexicographically optimal over the generated column set, within solver
tolerance, without a manually selected scalar resource price.

This formulation is conditional on the declared RCS scenario.  An uncertain-
RCS extension would replace $p_{qfb}$ by scenario coefficients
$p^{(s)}_{qfb}=p_{qfb}(\sigma_q^{(s)})$ and introduce scenario-wise deficits;
the expectation, worst-case, or CVaR aggregation must then be specified
explicitly.  Such a robust extension is not claimed in V1.2.

\subsection{Detector-Aligned Coarse-to-Fine Pricing}
Enumerating all bundles is combinatorial.  We first rank local and remote
observations separately using coarse delay--Doppler predictions, thereby
preventing either evidence family from disappearing from the candidate pool.
Generated bundles are then valued with the refined delay--Doppler table.  The
initial restricted master problem (RMP) contains an empty bundle, greedy
detector-marginal prefixes, and a singleton from each available local/remote
family for every admissible $(q,f)$ pair.

We relax bundle integrality and generate columns in lexicographic order.  The
worst-deficit LP is closed first; its optimum is then locked before pricing the
total-deficit LP.  Let $\beta_q$, $\pi_f^{\rm target}$,
$\pi_j^{\rm rx}$, $\pi_f^{\rm fusion}$, $\pi^{\rm report}$, and
$\pi^{\rm total}$ denote the corresponding nonnegative scarcity prices, and
let $\nu_q$ be the assignment dual.  The reduced cost of a candidate bundle is
\begin{align}
\bar c_{qfb}={}&-\nu_q-\beta_qp_{qfb}+\pi_f^{\rm target}
+\sum_j\pi_j^{\rm rx}a_{qfbj}\nonumber\\
&+\pi_f^{\rm fusion}c_{qfb}+\pi^{\rm report}r_{qfb}
+\pi^{\rm total}c_{qfb}.
\label{eq:v12-reduced-cost}
\end{align}
A bundle with $\bar c_{qfb}<0$ improves the current relaxed tier.  These
coefficients are dual variables of hard constraints rather than user-tuned
trade-off weights.

For a small shortlist, the pricing oracle enumerates every subset up to the
per-target cap.  For larger instances, it applies multi-start marginal pricing
from the empty, local-singleton, and remote-singleton states.  Exact pricing
certifies convergence only within the shortlisted column space.  Greedy
pricing is a scalable heuristic and carries no global-optimality claim.  After
pricing terminates, the accumulated RMP is solved as the exact four-tier mixed
integer problem.

\subsection{Validation Protocol}
Optimization correctness is evaluated for $M\in\{3,4\}$ and
$Q\in\{2,3\}$ by comparing all four lexicographic components against exhaustive
joint enumeration under identical local, receiver, fusion, assignment,
report, and total-observation constraints.  The full-scale study compares the
joint method with local-only bundles, fixed V1.1 fusion plus joint bundle
selection, V1.1 detector-aware sequential selection, and budgeted sensing-SINR
selection.  All methods share keyed primitive detector draws and report both
active-target $P_{FA}$ and its trial-cluster interval.  Large-scale performance
claims remain conditional on the preregistered independent Monte Carlo gate.
